God prepares a feast, yet some people will not come. Others come, yet their arrival is not the end of the story. Isaiah promises the destruction of death, while Paul still knows hunger. The Twenty-eighth Sunday's abundance is therefore more demanding than a picture of everyone enjoying plenty. It asks what kind of life belongs at God's table, and what trust in God's provision means before tears have been wiped away.

The first question comes into sharpest focus in the long Gospel. Gregory the Great identifies the wedding garment with charity: the guest has entered the Church but has not embraced the love proper to its Lord. Augustine's account of pardon and growth toward divine likeness explains why this demand neither purchases grace nor makes conversion unnecessary. The second question begins with Paul's acknowledgment of real want. Chrysostom explains why contentment and another person's gift belong together; Augustine tests the psalm's promise against the unsettling fact that a wicked person can remain wealthy while a faithful person is poor. Thomas Aquinas distinguishes the present banquet from its heavenly fullness.

These readings converge in Christ, who gives what his people cannot provide for themselves and draws them into a life of mutual care. Their emphases remain different. The first examines the heart of the guest already inside; the second examines the meaning of abundance while that guest can still be hungry. Mercy, illumination, offering and Communion give both questions their liturgical depth. The Church asks for help in doing good and for nourishment that changes those who receive it.

The Gospel may end with the gathering of the guests or continue through the judgment of the garment. Communion may be accompanied by the psalm of those who seek the Lord or by John's promise of likeness through vision. These alternatives make a real difference: the short Gospel does not proclaim the garment scene, and the Johannine Communion speaks explicitly of a fulfillment still to come. The complete set of appointments follows.
